\section{Harmonic Effects: Space Harmonics, Cogging, Crawling and Skew}

\subsection{Space Harmonics of 3-Phase MMF: Orders \texorpdfstring{$n = 6k \pm 1$}{n = 6k +/- 1}, Forward and Reverse Rotation}
In a balanced three-phase stator winding fed by sinusoidal currents, space harmonics of the magnetomotive force (MMF) are produced due to the discrete arrangement of coils in slots:
\begin{equation}
\boxed{n = 6k \pm 1 \qquad (k = 1, 2, 3, \dots)}
\end{equation}
\textbf{Properties of space harmonic orders:}
\begin{itemize}
  \item \textbf{Forward-rotating harmonics ($n = 6k+1$):} Orders $n = 7, 13, 19, 25, \dots$ rotate in the \textbf{same direction} as the fundamental at synchronous speed:
  \begin{equation}
  n_{s,n} = +\frac{n_s}{n} \quad [\mathrm{r/min}].
  \end{equation}
  \item \textbf{Reverse-rotating harmonics ($n = 6k-1$):} Orders $n = 5, 11, 17, 23, \dots$ rotate in the \textbf{opposite direction} to the fundamental at speed:
  \begin{equation}
  n_{s,n} = -\frac{n_s}{n} \quad [\mathrm{r/min}].
  \end{equation}
  \item \textbf{Triplen harmonics ($n = 3, 9, 15, \dots$):} In balanced three-phase systems, triplen MMF harmonics are displaced by $n\cdot 120^\circ = 360^\circ$ and sum algebraically to \textbf{zero} in the air gap.
  \item \textbf{Even harmonics ($n = 2, 4, 6, \dots$):} Cancel identically due to the half-wave anti-symmetry of the winding coils ($F(\theta+\pi) = -F(\theta)$).
\end{itemize}

\subsection{Cogging: Magnetic Locking at Standstill (\texorpdfstring{$S_1, S_2$}{S1, S2} Tooth Alignment)}
\textbf{Definition:} Cogging is the failure of an induction or permanent-magnet motor to start from standstill due to magnetic locking forces.

\textbf{Mechanism:}
\begin{itemize}
  \item Flux lines seek the path of minimum reluctance (acting like stretched elastic rubber bands).
  \item When stator slot count $S_1$ and rotor slot count $S_2$ are equal or have an integral ratio ($S_1/S_2 \in \mathbb{Z}$), all teeth align simultaneously at specific rotor positions.
  \item This creates large periodic detent alignment forces (cogging torque $T_{\mathrm{cog}}$) at zero speed.
  \item If $T_{\mathrm{cog}} > T_{\mathrm{starting}}$, the rotor locks magnetically and refuses to break away.
\end{itemize}
\textbf{Remedies:}
\begin{itemize}
  \item Skew rotor bars by one stator slot pitch ($\gamma_{\mathrm{sk}} = 2\pi p / S_1$).
  \item Select slot numbers such that $(S_1 - S_2)/(2p)$ is non-integral and $S_1/S_2$ is fractional.
  \item Use closed or semi-closed slots to reduce slot-opening permeance variations.
\end{itemize}

\subsection{Crawling: Asynchronous Torque Saddles (\texorpdfstring{$n_s/7, -n_s/5$}{ns/7, -ns/5}) and Synchronous Hooks}
\textbf{Definition:} Crawling is stable motor operation at a sub-synchronous speed (most commonly near $n_s/7$), preventing the motor from accelerating to its rated speed.

Crawling arises from two distinct physical mechanisms:

\textbf{1. Harmonic induction torques (asynchronous torque saddles):}
\begin{itemize}
  \item Stator space harmonics induce harmonic currents in the rotor cage, each producing an independent torque-speed characteristic.
  \item The \textbf{7th harmonic} rotates forward at synchronous speed $+n_s/7$. It produces positive motoring torque for $n_r < n_s/7$, passes through zero at $n_r = n_s/7$, and develops negative (generator/braking) torque for $n_r > n_s/7$.
  \item The combined fundamental + 7th torque curve develops a sharp dip or \textbf{saddle} around $n_r \approx n_s/7$.
  \item If the load torque curve intersects the rising portion of this saddle, the motor operating point locks stably at $n_r \approx n_s/7$.
  \item The \textbf{5th harmonic} rotates in reverse at synchronous speed $-n_s/5$, acting as a continuous brake opposing acceleration during startup.
\end{itemize}

\textbf{2. Harmonic synchronous torques (synchronous hooks):}
\begin{itemize}
  \item Occurs when a stator slot harmonic and a rotor slot harmonic have the \textbf{same number of poles} and rotate at the \textbf{exact same absolute speed} in the air gap.
  \item Stator slot harmonic order: $n = (S_1/p)\pm 1$ at speed $\pm n_s/n$.
  \item Rotor slot harmonic order: $n = (S_2/p)\pm 1$ at relative speed $\pm (n_s - n_r)/n$ with respect to the rotor.
  \item When they lock in step, they act like a miniature synchronous machine at that speed, producing a severe, sharp torque cusp (\textbf{hook}).
  \item \textbf{Example:} For $2p = 4$, $S_1 = 24$, $S_2 = 28$:
  Stator has 13th harmonic ($n_s/13$); rotor has 13th harmonic ($n_r - (n_s - n_r)/13$).
  In step when $n_s/13 = n_r - (n_s - n_r)/13 \implies n_r = n_s/7$.
  This produces a synchronous hook at $n_s/7$ driven by the \textbf{13th harmonics}, not the 7th!
\end{itemize}

\subsection{Tooth Harmonics Unattenuated by Winding Distribution: \texorpdfstring{$K_{dn} = K_{d1}$}{Kdn = Kd1}}
In a slotted $m$-phase stator winding with $q_1$ slots per pole per phase and $60^\circ$ phase spread ($\sigma = \pi/3$), slotting produces tooth harmonics at orders:
\begin{equation}
\boxed{n = 6k q_1 \pm 1 = \frac{S_1}{p} \pm 1 \qquad (k = 1, 2, \dots)}
\end{equation}
The EMF distribution factor for harmonic order $n$ is:
\begin{equation}
K_{dn} = \frac{\sin(n\sigma / 2)}{q_1 \sin(n\sigma / (2q_1))}.
\end{equation}
Substituting $n = 6k q_1 \pm 1$ and $\sigma = \pi/3$:
\begin{equation}
K_{dn} = \frac{\sin\left(\frac{(6k q_1 \pm 1)\pi}{6}\right)}{q_1 \sin\left(\frac{(6k q_1 \pm 1)\pi}{6q_1}\right)}
= \frac{\sin\left(k\pi \pm \frac{\pi}{6}\right)}{q_1 \sin\left(\frac{k\pi}{q_1} \pm \frac{\pi}{6q_1}\right)}
= \frac{\sin(\pi/6)}{q_1 \sin(\pi/(6q_1))} = K_{d1}.
\end{equation}
\textbf{Fundamental design conclusion:}
\begin{equation}
\boxed{K_{dn} = K_{d1} \quad \text{for all tooth harmonics } n = \frac{S_1}{p}\pm 1}
\end{equation}
Distributing coils among slots successfully attenuates low-order space harmonics (such as 5th and 7th), but \textbf{cannot attenuate tooth harmonics}. Skewing is therefore required!

\subsection{Skew Factor \texorpdfstring{$K_{\mathrm{sk}}$}{K(sk)} and Cancellation of Slot Harmonics (\texorpdfstring{$\gamma_{\mathrm{sk}} = 2\pi p / S_1$}{gamma(sk) = 2pi p / S1})}
Skewing the rotor bars by an electrical angle $\gamma_{\mathrm{sk}}$ spreads conductor segments over an arc, introducing a skew factor:
\begin{equation}
\boxed{K_{\mathrm{sk},1} = \frac{\sin(\gamma_{\mathrm{sk}}/2)}{\gamma_{\mathrm{sk}}/2}, \qquad K_{\mathrm{sk},n} = \frac{\sin(n\gamma_{\mathrm{sk}}/2)}{n\gamma_{\mathrm{sk}}/2}}
\end{equation}
where all angles inside the sine and denominator must be in \textbf{radians}.

\textbf{Skewing by exactly one stator slot pitch:}
The mechanical stator slot pitch is $\tau_{s1,\mathrm{mech}} = 2\pi / S_1$. In electrical radians:
\begin{equation}
\boxed{\gamma_{\mathrm{sk}} = p\,\tau_{s1,\mathrm{mech}} = \frac{2\pi p}{S_1} \quad [\mathrm{elect.\,rad}]}
\end{equation}
Evaluating the skew factor for stator slot harmonics of order $n = \frac{S_1}{p}\pm 1$:
\begin{equation}
\frac{n\gamma_{\mathrm{sk}}}{2} = \frac{1}{2}\left(\frac{S_1}{p}\pm 1\right)\left(\frac{2\pi p}{S_1}\right) = \pi \pm \frac{\pi p}{S_1}.
\end{equation}
Taking the sine:
\begin{equation}
\sin\left(\frac{n\gamma_{\mathrm{sk}}}{2}\right) = \sin\left(\pi \pm \frac{\pi p}{S_1}\right) = \mp\sin\left(\frac{\pi p}{S_1}\right).
\end{equation}
Since slot count per pole-pair $S_1/p \gg 1$, the angle $\frac{\pi p}{S_1}$ is small, yielding:
\begin{equation}
\boxed{K_{\mathrm{sk},n} \approx 0 \quad \text{for } n = \frac{S_1}{p}\pm 1}
\end{equation}
\textbf{Physical reason:} Skewing by one stator slot pitch establishes an effective phase spread of $2\pi\,\mathrm{rad}$ (one full wave cycle) along the axial core length at the slot harmonic wavelength. The positive and negative half-waves induced along the bar cancel out completely.
The fundamental is only slightly diminished: $K_{\mathrm{sk},1} \approx 0.98\text{--}0.99$.

\subsection{Summary Table: Slot Combination Rules to Prevent Parasitic Effects}
\begin{center}
\begin{tabular}{@{}p{0.24\textwidth}p{0.32\textwidth}p{0.40\textwidth}@{}}
\toprule
\textbf{Parasitic Phenomenon} & \textbf{Harmful Condition} & \textbf{Design Rule to Eliminate / Avoid} \\
\midrule
\textbf{Cogging} & $S_1 = S_2$ or $S_1/S_2 \in \mathbb{Z}$ & $(S_1 - S_2)/(2p) \neq \text{integer}$, $S_1/S_2$ fractional, skew slots \\
\textbf{Harmonic induction torque} & Large rotor bar count & $S_2 < 1.25(S_1 + p)$ \\
\textbf{Harmonic synchronous torque} & Synchronous speed lock & $S_1 \neq S_2 \quad\text{and}\quad S_1 - S_2 \neq \pm 2p$ \\
\textbf{Unbalanced Magnetic Pull} & Pole count difference $= 2$ & $S_1 - S_2 \neq \pm 1 \quad\text{and}\quad S_1 - S_2 \neq \pm 2p \pm 1$ \\
\bottomrule
\end{tabular}
\end{center}

\subsection{Exam Recipe: Calculating Skew Factors and Harmonic EMF Reduction}
\textbf{Use when:} Problem specifies stator slot count $S_1$, pole pairs $p$ (or poles $2p$), and asks for skew distribution factors $K_{\mathrm{sk},1}, K_{\mathrm{sk},n}$ or reduced harmonic voltages when rotor is skewed by one stator slot pitch.

\textbf{Steps:}
\begin{enumerate}
  \item \textbf{Determine slots per pole:}
  \[
  \text{Slots per pole} = \frac{S_1}{2p}.
  \]
  \item \textbf{Calculate electrical slot pitch angle $\gamma_{\mathrm{sk}}$:}
  \[
  \gamma_{\mathrm{sk}} = \frac{2\pi p}{S_1} = \frac{\pi}{S_1/(2p)} \quad [\mathrm{rad}].
  \]
  \item \textbf{Compute fundamental skew factor:}
  \[
  K_{\mathrm{sk},1} = \frac{\sin(\gamma_{\mathrm{sk}}/2)}{\gamma_{\mathrm{sk}}/2}.
  \]
  \item \textbf{Identify tooth harmonic orders:}
  \[
  n = \frac{S_1}{p} \pm 1.
  \]
  \item \textbf{Compute harmonic skew factor:}
  \[
  K_{\mathrm{sk},n} = \frac{\sin(n\gamma_{\mathrm{sk}}/2)}{n\gamma_{\mathrm{sk}}/2}.
  \]
  \item \textbf{Compute skewed induced voltages:}
  \[
  E_{1,\mathrm{skewed}} = E_{1,\mathrm{unskewed}}\cdot K_{\mathrm{sk},1}, \qquad E_{n,\mathrm{skewed}} = E_{n,\mathrm{unskewed}}\cdot |K_{\mathrm{sk},n}|.
  \]
  \item \textbf{Conclusion/Comment:} Note that $K_{\mathrm{sk},n} \ll 1$ (e.g., $< 0.05$), confirming that skewing by one stator slot pitch suppresses slot harmonics and eliminates cogging torque.
\end{enumerate}

\textbf{Fast checks:}
\begin{itemize}
  \item Fundamental factor $K_{\mathrm{sk},1}$ is always very close to 1 (typically $0.98\text{--}0.999$).
  \item Harmonic factor $|K_{\mathrm{sk},n}|$ for $n = S_1/p \pm 1$ is always very small (typically $< 0.1$).
\end{itemize}

\textbf{Common traps:}
\begin{itemize}
  \item Using degrees instead of radians in the denominator of $\frac{\sin(\theta)}{\theta}$.
  \item Using mechanical slot pitch $\tau_{s1} = 2\pi / S_1$ instead of electrical slot angle $\gamma_{\mathrm{sk}} = 2\pi p / S_1$.
  \item Forgetting that $K_{dn} = K_{d1}$ for tooth harmonics (winding distribution cannot eliminate them; only skewing does).
\end{itemize}
